\chapter{Elderly Mental Health}
\index{Elderly Mental Health}

\section*{Definition and Overview}
Mental health in later life is just as important as physical health, yet mental health conditions in older people are often missed, dismissed, or attributed to ageing. Depression, anxiety, and dementia are the most common mental health problems in older adults, and they frequently coexist with physical illness, disability, and social isolation. Mental illness is not a normal part of ageing, and most conditions are treatable. This chapter explains the mental health conditions that affect older people, the barriers to recognition, and the treatments and support available.

\section*{Epidemiology}
Depression affects a significant minority of older people living in the community, and higher proportions of those in hospital and residential care. Anxiety disorders are also common, and dementia affects a growing number of people as the population ages. Older men have the highest rate of suicide of any age group, yet mental health problems in older people are under-treated, partly because symptoms are mistaken for ageing and partly because older people are less likely to seek help.

\section*{Aetiology and Pathophysiology}
Mental health problems in later life arise from a combination of biological, psychological, and social factors.
\begin{itemize}
    \item \textbf{Physical illness:} chronic disease, pain, and disability increase the risk of depression and anxiety
    \item \textbf{Brain changes:} vascular disease, stroke, and neurodegenerative changes contribute to depression and dementia
    \item \textbf{Medication:} many medicines used by older people affect mood, sleep, and cognition
    \item \textbf{Loss and grief:} bereavement, retirement, and loss of independence are common in later life
    \item \textbf{Social isolation:} loneliness and reduced social contact are powerful risk factors
    \item \textbf{Caregiving:} caring for a partner with illness or disability increases stress and depression risk
    \item \textbf{Stigma:} older people are less likely to seek help because of the belief that mental illness is a personal failing
\end{itemize}

\section*{Clinical Features}
Mental health problems in older people may present differently from in younger people.
\begin{itemize}
    \item Low mood, loss of interest, and withdrawal from activities
    \item Irritability, agitation, and anxiety
    \item Sleep disturbance and loss of appetite
    \item Unexplained physical symptoms, such as pain, fatigue, and dizziness
    \item Memory problems and difficulty concentrating
    \item Slowing of thought and movement
    \item Neglect of personal care and medications
    \item Confusion that can be mistaken for dementia
    \item Thoughts of death or suicide
\end{itemize}

\section*{Differential Diagnosis}
Symptoms that suggest mental illness can have other causes.
\begin{itemize}
    \item Dementia, including Alzheimer's disease and vascular dementia
    \item Delirium, an acute confusional state with an underlying medical cause
    \item Medication side effects, including from sedatives and corticosteroids
    \item Thyroid disease, anaemia, and vitamin deficiencies
    \item The normal grieving process after bereavement
    \item The effects of alcohol and other substances
\end{itemize}

\section*{When to Seek Medical Care}
Families and carers should encourage older people to see a GP for mood, memory, sleep, or anxiety concerns, and the GP should routinely ask about mental health. Changes that are new, persistent, or affecting function always warrant assessment.

\textbf{Contact your local emergency services for:}
\begin{itemize}
    \item Thoughts of self-harm or suicide, or talking about wanting to die
    \item Sudden confusion, which may indicate a medical emergency
    \item Severe agitation, paranoia, or hallucinations
    \item Collapse, stroke symptoms, or signs of a medical emergency
\end{itemize}

\section*{Diagnosis and Investigations}
Assessment distinguishes mental illness from other causes of the symptoms.
\begin{itemize}
    \item \textbf{History:} symptoms, their onset, medicines, physical health, and social circumstances, gathered from the person and family
    \item \textbf{Mental state examination:} mood, cognition, and risk assessment
    \item \textbf{Cognitive testing:} brief tests distinguish depression-related memory problems from dementia
    \item \textbf{Blood tests:} thyroid function, B12, folate, and other checks
    \item \textbf{Medication review:} identifying medicines that affect mood or cognition
    \item \textbf{Referral:} to aged psychiatry, psychology, or dementia services as needed
\end{itemize}

\section*{Management}
Management combines treatment of the condition with support for daily life.
\begin{itemize}
    \item \textbf{Psychological therapy:} cognitive behavioural therapy and counselling are effective in older people
    \item \textbf{Antidepressants:} effective for moderate to severe depression, chosen with attention to side effects
    \item \textbf{Treating the physical:} management of pain, illness, and medication effects
    \item \textbf{Social connection:} groups, volunteering, and community activities reduce isolation
    \item \textbf{Exercise:} physical activity improves mood, sleep, and cognition
    \item \textbf{Supporting carers:} respite care and carer education
    \item \textbf{Specialist care:} aged psychiatry for complex or severe illness
    \item \textbf{Suicide prevention:} safety planning and crisis support where risk exists
\end{itemize}

\section*{Complications}
\begin{itemize}
    \item Worsening physical health and disability
    \item Falls and reduced independence
    \item Social withdrawal and isolation
    \item Poor adherence to treatment for physical illness
    \item Premature admission to residential care
    \item Suicide, particularly in older men
    \item Carer stress and burnout
\end{itemize}

\section*{Prognosis and Outcomes}
Mental health conditions in later life respond well to treatment, and most older people with depression improve significantly with appropriate care. Outcomes depend on recognition, treatment, and social support. Ageing psychiatry and psychology services are effective, and many older people achieve remission and regain a good quality of life.

\section*{Prevention and Self-Care}
\begin{itemize}
    \item Stay socially connected through family, friends, and community groups
    \item Keep physically active and eat well
    \item Maintain a regular sleep routine
    \item Review medicines with a doctor or pharmacist regularly
    \item Speak to a GP about mood or memory concerns early
    \item Plan meaningful activities and interests after retirement
    \item Seek help for grief and loss rather than struggling alone
\end{itemize}

\section*{Sources and Further Reading}
The following reputable sources provide further information for healthcare students and patients seeking reliable, current advice about mental health in older people.
\begin{itemize}
    \item Beyond Blue. \textit{Depression and anxiety in older people.} https://www.beyondblue.org.au/
    \item Healthdirect Australia. \textit{Mental health in older people.} https://www.healthdirect.gov.au/
    \item NHS. \textit{Mental health for older adults.} https://www.nhs.uk/mental-health/
    \item Older Persons Mental Health services. \textit{Aged psychiatry information.} https://www.health.nsw.gov.au/
    \item Dementia Australia. \textit{Memory and dementia support.} https://www.dementia.org.au/
    \item Lifeline. \textit{Crisis support.} https://www.lifeline.org.au/
\end{itemize}
